\section{The Cubic--Quintic Scheme}
\label{sec:scheme}

The previous section lets us work directly with fixed-dimensional limiting
Krivine schemes. We now define the two-dimensional scheme used for our upper
bound. Its first coordinate combines the correlation powers \(t,t^3,t^5\),
while its second coordinate keeps correlation \(t\). In the next two sections,
we compute its correlation function and certify its inverse majorant.

We first describe the general form of the construction. Let
\(D\subseteq\{3,5,7,\ldots\}\) be finite, choose weights
\((s_d)_{d\in D}\), and set
\[
    \sigma_d=(-1)^{(d-1)/2},
    \qquad
    V_D=1+\sum_{d\in D}s_d^2.
\]
A degree-\(d\) Hermite direction contributes the correlation \(t^d\).
After combining these directions into one normalized Gaussian coordinate,
the resulting correlation map is
\begin{equation}
    \rho_D(t)
    =
    \frac{t+\sum_{d\in D}\sigma_d s_d^2 t^d}{V_D}.
    \label{eq:general-hermite-correlation-map}
\end{equation}
The map \(\rho_D\) is allowable, since the absolute values of its
coefficients sum to one. Thus, for any \(\eta\in\mathbb{R}\), the partitions
\[
    \operatorname{sgn}\left(
        w+\frac{\eta}{\sqrt{V_D}}\He_3(x)
    \right)
    \qquad\text{and}\qquad
    \operatorname{sgn}\left(
        w-\frac{\eta}{\sqrt{V_D}}\He_3(x)
    \right)
\]
together with the correlation maps \(\rho_D(t)\) and \(t\) define a
two-dimensional limiting Krivine scheme.

For the scheme used in our proof, we take \(D=\{3,5\}\). The cubic
direction is anti-aligned between the two sides, while the quintic direction
is aligned. We use the parameters
\begin{equation}
    \eta=0.136419125,
    \qquad
    s_3=0.34101124,
    \qquad
    s_5=0.05276111.
    \label{eq:params}
\end{equation}
Set
\[
    V=1+s_3^2+s_5^2,
    \qquad
    \vartheta=\frac{\eta}{\sqrt V},
\]
and define
\begin{equation}
    \rho(t)
    =
    \frac{t-s_3^2t^3+s_5^2t^5}{V}.
    \label{eq:rho-collapsed}
\end{equation}

\begin{definition}[Cubic--quintic limiting scheme]
\label{def:cubic-quintic-scheme}
Define \(f,g:\mathbb{R}^2\to\{-1,1\}\) by
\begin{equation}
    f(w,x)
    =
    \sgn\bigl(w+\vartheta\He_3(x)\bigr),
    \qquad
    g(w,x)
    =
    \sgn\bigl(w-\vartheta\He_3(x)\bigr).
    \label{eq:cubic-quintic-partitions}
\end{equation}
Our construction is the limiting Krivine scheme
\[
    \mathcal{S}=(f,g,\rho_1,\rho_2),
    \qquad
    \rho_1(t)=\rho(t),
    \qquad
    \rho_2(t)=t.
\]
We write \(H\) for its correlation function.
\end{definition}

The map \(\rho\) is allowable because the absolute values of its
coefficients sum to \((1+s_3^2+s_5^2)/V=1\). The partitions are odd and
satisfy \(g(w,x)=f(w,-x)\). Their threshold boundaries also have Gaussian
measure zero, so the approximation theorem from the previous section
applies.

More explicitly, let \((W,X)\) and \((W',X')\) be standard Gaussian vectors.
Assume that the two coordinate pairs are independent and that
\(\operatorname{corr}(W,W')=\rho(t)\) and
\(\operatorname{corr}(X,X')=t\). Then
\begin{equation}
    H(t)
    =
    \frac{\pi}{2}
    \E\bigl[f(W,X)g(W',X')\bigr].
    \label{eq:lim-H}
\end{equation}

The signs and weights in \(\rho\) are chosen to make the cubic and quintic
coefficients of \(H\) nearly vanish. This in turn suppresses the first two
nonlinear coefficients of \(H^{-1}\). Section~\ref{sec:coefficients} makes
this relation explicit.

\begin{theorem}[Cubic--quintic upper bound]
\label{thm:cubic-quintic-upper-bound}
Let \(M\) be the inverse majorant of \(H\). At
\(\gamma=0.881545409\), we have \(M(\gamma)<1\). Consequently,
\[
    K_G
    \leq
    \frac{\pi}{2\gamma}
    \leq
    1.7818666069360661
    <
    \frac{\pi}{2\log(1+\sqrt{2})}
    -
    3.4737\cdot 10^{-4}.
\]
\end{theorem}